%!TEX program = pdflatex
\documentclass[aps,%
onecolumn,%
floatfix,%
pra,%
amsfonts,%
longbibliography,%
groupedaddress,%
superscriptaddress]{revtex4-2}%

%%%%%%%%%%%%%%%%%%%%%%%%%%%%%%%%%%%%%%%%%%%%%%%%%%%%%%%%%%%%%%%%%%%%%%%%%
% Packages
%%%%%%%%%%%%%%%%%%%%%%%%%%%%%%%%%%%%%%%%%%%%%%%%%%%%%%%%%%%%%%%%%%%%%%%%%
\usepackage{amsmath}
\usepackage{amssymb}
\usepackage{amsthm}
\usepackage{bm,bbm} % to use the command \mathbbm{1}, \bm{x}
\usepackage{graphicx}
\usepackage{indentfirst}
\usepackage{enumitem}
% RevTeX provides natbib; do not load biblatex here.
% \usepackage{algorithmicx}
% \usepackage{algpseudocode}
% \usepackage{algorithm}
% \usepackage{algorithmic}
% \usepackage{algcompatible}
\usepackage{algpseudocode}
\usepackage{tikz}
\usepackage{quantikz}
\usetikzlibrary{graphs} 
\usepackage[titletoc,toc,page,title]{appendix}
\definecolor{beamer@blendedblue}{rgb}{0.2,0.2,0.7}      % color used in beamer
\usepackage[bookmarks=false]{hyperref}
\hypersetup{colorlinks=true,%
citecolor=beamer@blendedblue,%
linkcolor=beamer@blendedblue,%
urlcolor=beamer@blendedblue,%
pdfstartview=FitH,%
bookmarksopen=true}

\def\UrlBreaks{\do\/\do-} % url break
\usepackage[T1]{fontenc}
\usepackage{times}
\usepackage{booktabs} % to use the \toprule \midrule command
%%%%%%%%%%%%%%%%%%%%%%%%%%%%%%%%%%%%%%%%%%%%%%%%%%%%%%%%%%%%%%%%%%%%%%%%%
% General Commands
%%%%%%%%%%%%%%%%%%%%%%%%%%%%%%%%%%%%%%%%%%%%%%%%%%%%%%%%%%%%%%%%%%%%%%%%%
\newtheorem{theorem}{Theorem}
\newtheorem{proposition}[theorem]{Proposition}
\newtheorem{lemma}{Lemma}
\newtheorem{corollary}{Corollary}

%%%%%%%%%%%%%%%%%%%%%%%%%%%%%%%%%%%%%%%%%%%%%%%%%%%%%%%%%%%%%%%%%%%%%%%%%
\usepackage{pifont}

%%%%%%%%%%%%%%%%%%%%%%%%%%%%%%%%%%%%%%%%%%%%%%%%%%%%%%%%%%%%%%%%%%%%%%%%%
% Quantum Commands
%%%%%%%%%%%%%%%%%%%%%%%%%%%%%%%%%%%%%%%%%%%%%%%%%%%%%%%%%%%%%%%%%%%%%%%%%
\DeclareMathOperator{\Tr}{Tr}
\DeclareMathOperator{\tr}{Tr}
\DeclareMathOperator{\NOT}{NOT}
\DeclareMathOperator{\diag}{diag}
\DeclareMathOperator{\rank}{rank}
\DeclareMathOperator{\conv}{conv}
\DeclareMathOperator{\sgn}{sgn}
\DeclareMathOperator{\Sp}{Span}
\DeclareMathOperator{\id}{id}
% \newcommand{\bra}[1]{\langle #1\rvert}
% \newcommand{\ket}[1]{\lvert #1\rangle}
% \newcommand{\proj}[1]{\lvert #1\rangle\!\langle #1\rvert}
\newcommand{\ox}{\otimes}
\newcommand{\1}{\mathbbm{1}}
\newcommand{\mean}[1]{\langle #1\rangle}
\newcommand{\gen}[1]{\langle #1\rangle}
\newcommand{\ketbra}[2]{\lvert #1\rangle\!\langle #2\rvert}
% \newcommand{\braket}[2]{\langle #1\vert #2\rangle}
%%%%%%%%%%%%%%%%%%%%%%%%%%%%%%%%%%%%%%%%%%%%%%%%%%%%%%%%%%%%%%%%%%%%%%%%%
%%% Mathematical symbols
%%%%%%%%%%%%%%%%%%%%%%%%%%%%%%%%%%%%%%%%%%%%%%%%%%%%%%%%%%%%%%%%%%%%%%%%%
\newcommand*{\cF}{\mathcal{F}}
\newcommand*{\cH}{\mathcal{H}}
\newcommand*{\cO}{\mathcal{O}}
\newcommand*{\cP}{\mathcal{P}}
\newcommand*{\cS}{\mathcal{S}}
\newcommand*{\bF}{\mathbb{F}}
\newcommand*{\bP}{\mathbb{P}}

% Colors
\definecolor{alizarin}{rgb}{0.82, 0.1, 0.26}
\definecolor{googleblue}{HTML}{4285F4}
\definecolor{googlered}{HTML}{DB4437}
\definecolor{googleyellow}{HTML}{F4B400}
\definecolor{googlegreen}{HTML}{0F9D58}
\definecolor{klevinblue}{HTML}{002FA7}
\definecolor{tiffanyblue}{HTML}{0ABAB5}

\newcommand{\KW}[1]{{\color{googlered}(KW: #1)}}
\newcommand{\SYC}[1]{\textcolor{blue}{(SYC: #1)}}

\begin{document}

\title{Simulation of Local addressing protocol}
%%%
\author{Siyuan Chen}
\affiliation{Hefei National Research Center for Physical Sciences at the Microscale and School of Physical Sciences, 
University of Science and Technology of China, Hefei 230026, China}%
%%%


\begin{abstract}
\noindent 
In this work, we simulate the local addressing protocol used by Evered \emph{et.al.}~\cite{Evered_2023}, based on $^{87}Rb$ atoms.
\end{abstract}
\date{\today}
\maketitle

\section{Results}

\subsection{ground state shift v.s. Scattering rate}
Local addressing laser is near-resonant between transitions~\cite{steck2001rubidium}:
\begin{align*}
    &5S_{1/2} \leftrightarrow 5P_{3/2}, \quad \lambda = 780 \mathrm{nm}, \quad \Gamma = 2\pi \times 6.065(9) \mathrm{MHz},\\
    &5S_{1/2} \leftrightarrow 5P_{1/2}, \quad \lambda = 794.8 \mathrm{nm}, \quad \Gamma = 2\pi \times 5.746(8) \mathrm{MHz}.
\end{align*}
For each near-resonant transition, the dipole potential and the scattering rate in the relevant case of large detunings and negligible saturation, we have the following expressions~\cite{grimm1999opticaldipoletrapsneutral}:
\begin{align*}
    \Delta_{LS}(\omega)
&=  - \frac{3 \pi c^2I}{2 \omega_0^3} \,
\left( \frac{\Gamma}{\omega_0 -\omega} + 
\frac{\Gamma}{\omega_0 +\omega} \right) \propto  - \frac{1}{\omega_0^3} \,
\left( \frac{\Gamma}{\omega_0 -\omega} + 
\frac{\Gamma}{\omega_0 +\omega} \right)  , \\
\Gamma_{\rm sc}(\omega) &=  \frac{3 \pi c^2I }{2 \hbar \omega_0^3} \,
\left(\frac{\omega}{\omega_0}\right)^3 \,
\left( \frac{\Gamma}{\omega_0 -\omega} + 
\frac{\Gamma}{\omega_0 +\omega} \right)^2  \propto  - \frac{1}{\omega_0^3} \left(\frac{\omega}{\omega_0}\right)^3\,
\left( \frac{\Gamma}{\omega_0 -\omega} + 
\frac{\Gamma}{\omega_0 +\omega} \right)^2  .
\end{align*}
Since $5P_{3/2}$ and $5P_{1/2}$ each contains $4$ and $2$ degenerate states, respectively, we reweighted the lightshift and the scattering rate by fraction $U_{\rm total}(\omega) = \frac{2}{3} U^{\rm 780}(\omega) + \frac{1}{3} U^{\rm 794}(\omega)$ and $\Gamma_{\rm total}(\omega) = \frac{2}{3} \Gamma^{\rm 780}(\omega) + \frac{1}{3} \Gamma^{\rm 794}(\omega)$. To increase the accuracy, instead of calculating the overall constant ${3 \pi c^2I }/{2 \hbar \omega_0^3}$ directly, we rescaled to match the measured differential AC Stark shift in \cite{manovitz2025quantum} with $\Delta_{LS} = 12.2(3) \mathrm{MHz}, \Gamma_{sc} = 35 \mathrm{Hz}$ at laser power $P \approx 160 \mathrm{\mu W}$ and laser wavelength $\lambda = 784 \mathrm{nm}$.

\begin{figure}[h]
\centering
\begin{minipage}[t]{0.4\textwidth}
  \centering
  \includegraphics[width=\linewidth]{Figures/RbProperties.png}
\end{minipage}
\begin{minipage}[t]{0.55\textwidth}
  \centering
  \includegraphics[width=\linewidth]{Figures/ac_stark_profiles.png}
\end{minipage}
\caption{Left: Energy level of $^{87}Rb$ that near the laser frequency of 780 nm. Right: Ac Stark landscape of $^{87}Rb$ that near the laser frequency of 780 nm. }
\end{figure}

\subsection{Behavior under adiabatic evolution}

In the Rydberg experiment~\cite{manovitz2025quantum}, the local addressing protocol is used to address certain atoms to the ground state while the other atoms are adiabatically evolved under Rydberg interaction and driving of $420 \mathrm{nm}$ and $1013 \mathrm{nm}$ lasers. 
A 2-atom version of the Hamiltonian, each labeled by $i =1,2$, is specified by operators $X_i = \ket{g_i}\bra{r_i} + \ket{r_i}\bra{g_i}$ and $n_i = \ket{r_i}\bra{r_i}$ and $V_{12} = C_6/R^6$ with $R$ the distance between the two atoms:
\begin{align*}
    H &= \frac{\Omega_{\rm eff}(t)}{2} \left(X_1 + X_2 \right) - \left( \Delta(t) + \delta_1(\omega) \right) n_1 - \left( \Delta(t) + \delta_2(\omega) \right) n_2 + V_{12} n_1 n_2.
\end{align*}
Atoms are put at distance $R = 4 \mu m$, which cause the Rydberg interaction intensity $V_{12} = C/R^6 = 213.4 \mathrm{MHz}$.

\paragraph*{The control of $\Omega_{\rm eff}(t)$.}
The $\Omega_{\rm eff}$ is the effective Rabi frequency of two-photon transition process driven by the 420nm and 1013nm lasers along:
\begin{align*}
    \ket{5S_{1/2}, F = 2, m = -2} \xrightarrow[\sigma_{-}]{\Omega_{\rm 420}(t) e^{i \varphi(t)}} \ket{6P_{3/2},\, m_J = -3/2,\, m_I = -3/2} \xrightarrow[\sigma_{+}]{\Omega_{\rm 1013}} \ket{70S_{1/2},\, m_J = -1/2,\, m_I = -3/2}.
\end{align*}
Refer to~\cite{manovitz2025quantum}, we set the 420nm and 1013nm Rabi frequency both to be $135 \mathrm{MHz}$ and blue detuning of the mid-state $6P_{3/2}$ with $\Delta = 2.4 \mathrm{GHZ}$.
This introduces no effective detuning and the effective Rabi frequency $\Omega_{\text{eff}} = {\Omega_{420} \Omega_{1013}}/{2\Delta} = 3.8 \mathrm{MHz}$ that meets the setting in the reference~\cite{manovitz2025quantum}.
To control $\Omega_{\rm eff}$ changing with time, we ramp up the Rabi frequency of the 420nm laser.
In the reference~\cite{manovitz2025quantum}, it ramp-up the Rabi-frequency of the 420nm laser without specifying the specific ramping function.
In our simulation, we ramp up the 420-Rabi with Blackman pulse, $A(t)$.
In the case $\Delta(0) = -\infty$, $H \approx n_1 + 20n_2 + 10 V_{12} n_1 n_2 $ and take $\ket{gg}$ as ground state. In the case $\Delta(0) = \infty$, depending on the local addressing laser, the ground state is $\ket{gr}$ or $\ket{rg}$.


\paragraph*{The control of $\Delta(t)$.}
The phase $\varphi(t)$ of the 420nm laser is a quadratic function encoding a linear chirp:
$$
\varphi(t) = \omega_0\,t + \frac{\alpha}{2}\,t^2,\quad \dot{\varphi}(t) = \omega_0 + \alpha\,t = \delta_{\rm start} + (\delta_{\rm end} - \delta_{\rm start})\,\frac{t}{T}.
$$
We setup $\delta_{\rm start} = -40 \mathrm{MHz}$ and $\delta_{\rm end} = 40 \mathrm{MHz}$, swept within the time $T = 4.5\mu s$.
With the Blackman pulse $A(t)$, an effective detuning is added to the detuning $\dot{\varphi}(t)$ as $\Delta_{\rm eff}(t) = A^2(t)/4 \Delta(t)$.

\begin{figure}[h]
    \centering
    \begin{minipage}[t]{0.6\textwidth}
      \centering
      \includegraphics[width=\textwidth]{Figures/addressing_enlevel.png}
    \end{minipage}
    \begin{minipage}[t]{0.38\textwidth}
      \centering
      \includegraphics[width=\textwidth]{Figures/fig4_420_phase_diagnostics.png}
    \end{minipage}
    \caption{Left: adiabatic evolution of the ground state shift and the scattering rate. Right: phase and amplitude diagnostics of the 420nm laser. }
\end{figure}

\paragraph*{The control of $\delta_i(\omega)$.}
Along with these two sweepping lasers, we also set up a local addressing laser, with wavelength rangeing from 781 nm to 786 nm, and laser power ranging from 140mW to 500 mW.
This local addressing laser operates on atom 1 only, while its Gaussian shape introduces a cross-talk to atom 2 with a reduction $e^{-2(R/w)^2}$ proportionally.
Effectively, we take both $w = 1 \mu m$ and $w = 2 \mu m$ to represent the local addressing laser.
 

\section{Results and conclusion}

\begin{figure}[h]
    \centering
    \begin{minipage}[t]{0.3\textwidth}
      \centering
      \includegraphics[width=\textwidth]{Figures/exp1_addressing_opt_smoothed.png}
    \end{minipage}
    \begin{minipage}[t]{0.3\textwidth}
      \centering
      \includegraphics[width=\textwidth]{Figures/exp2_addressing_opt_smoothed.png}
    \end{minipage}
    \begin{minipage}[t]{0.3\textwidth}
        \centering
        \includegraphics[width=\textwidth]{Figures/exp3_addressing_opt_smoothed.png}
      \end{minipage}

      \includegraphics[width=\textwidth]{Figures/exp3_addressing_opt_components.png}
    \caption{Representative optimization landscapes. Top row: Gaussian-smoothed total-cost maps for exp1--exp3. Bottom: cost decomposition for exp3. The same optimization was also carried out for exp4 and exp5; their optimal operating points are summarized below.}
\end{figure}

We compare five scans of the local addressing beam by varying the wavelength
$\lambda \in [781,786]~\mathrm{nm}$ and the power
$P \in [140,500]~\mu\mathrm{W}$. For each scan, we minimize the cost
function
\begin{align*}
    E_{\rm total} = E_{\rm leak} + E_{\rm xtalk} + E_{\rm sc},
\end{align*}
where $E_{\rm leak}=P_{rg}+P_{rr}$ measures imperfect pinning,
$E_{\rm xtalk}=P_{gg}$ measures residual population on the spectator atom, and
$E_{\rm sc}=1-e^{-\Gamma_{\rm sc}T}$ estimates the scattering penalty during
the sweep.

\begin{table}[h]
    \centering
    \caption{Best operating point found in each wavelength--power scan.}
    \label{tab:addressing_summary}
    \begin{tabular}{lccccccc}
        \toprule
        Scan & $R$ ($\mu$m) & $w$ ($\mu$m) & $T$ ($\mu$s) & $\Delta_{\max}$ (MHz) & $\lambda_{\rm opt}$ (nm) & $P_{\rm opt}$ ($\mu$W) & $E_{\rm total}^{\rm min}$ \\
        \midrule
        exp1 & 4.0 & 1.0 & 4.5 & 40.0 & 785.58 & 204.6 & $2.55\times10^{-4}$ \\
        exp2 & 4.0 & 1.0 & 9.0 & 40.0 & 786.00 & 204.6 & $5.90\times10^{-4}$ \\
        exp3 & 4.5 & 2.0 & 4.5 & 40.0 & 785.58 & 204.6 & $2.77\times10^{-4}$ \\
        exp4 & 4.5 & 2.0 & 4.5 & 11.4 & 785.58 & 490.8 & $3.62\times10^{-3}$ \\
        exp5 & 4.5 & 2.0 & 9.0 & 80.0 & 785.24 & 444.6 & $9.79\times10^{-4}$ \\
        \bottomrule
    \end{tabular}
\end{table}

\subsection{Comparison of the five scans}

\begin{enumerate}[label=(\roman*)]
    \item The optimal wavelength is remarkably stable across all five scans.
    The minimum always appears in a narrow window between $785.24$ and
    $786.00~\mathrm{nm}$. This is slightly red-detuned from the experimental
    calibration point at $784~\mathrm{nm}$, where the measured differential AC
    Stark shift is about $12.2~\mathrm{MHz}$ at
    $160~\mu\mathrm{W}$ with a scattering rate of about $35~\mathrm{Hz}$.
    Therefore, the simulation reproduces the experimental calibration at
    $784~\mathrm{nm}$, but the optimization consistently prefers a slightly
    larger wavelength.

    \item The optimal power is not universal; it depends strongly on the sweep
    protocol. In exp1--exp3, where the detuning window is
    $\pm 40~\mathrm{MHz}$, the optimum is always close to
    $205~\mu\mathrm{W}$ and gives a local shift
    $\delta_A \approx -8$ to $-10~\mathrm{MHz}$. By contrast, in exp5, where
    the sweep window is enlarged to $\pm 80~\mathrm{MHz}$ and the gate time is
    extended to $9~\mu\mathrm{s}$, the optimum moves to
    $444.6~\mu\mathrm{W}$ and gives
    $\delta_A \approx -22.9~\mathrm{MHz}$. Hence the laser power should be
    optimized together with the chirp window and the gate time, rather than
    quoted as a stand-alone parameter.

    \item The geometry change from exp1 to exp3 has only a small effect on the
    optimum. When the atom spacing and beam waist are changed from
    $(R,w)=(4~\mu\mathrm{m},1~\mu\mathrm{m})$ to
    $(4.5~\mu\mathrm{m},2~\mu\mathrm{m})$, the optimizer still selects the
    same point, $\lambda_{\rm opt}=785.58~\mathrm{nm}$ and
    $P_{\rm opt}=204.6~\mu\mathrm{W}$. The minimum cost only changes from
    $2.55\times 10^{-4}$ to $2.77\times 10^{-4}$. This suggests that, for the
    present two-atom geometry, the optimum is mainly determined by the balance
    between pinning quality and scattering, while the residual Gaussian tail on
    the spectator atom is already small enough not to move the optimum
    appreciably.

    \item Exp4 shows that insufficient sweep range cannot be compensated for by
    simply increasing the power. In this scan, the detuning window is reduced
    to only $\pm 11.4~\mathrm{MHz}$. The optimizer then pushes the power to
    $490.8~\mu\mathrm{W}$, but the minimum cost still rises to
    $3.62\times 10^{-3}$, which is about one order of magnitude worse than in
    exp1, exp3, or exp5. The main error in this regime is crosstalk, which
    means that a too narrow chirp window does not provide enough adiabatic
    separation even when the local light shift is increased.

    \item Among the five scans, exp5 gives the best strong-pinning protocol.
    Its optimum is located at
    $(\lambda,P)=(785.24~\mathrm{nm},444.6~\mu\mathrm{W})$, where
    $P_{gr}=0.9988$, the crosstalk error is reduced to
    $2\times 10^{-6}$, the leakage error is
    $4.47\times 10^{-4}$, and the scattering penalty is
    $5.30\times 10^{-4}$. In other words, once the sweep window is sufficiently
    wide, crosstalk is nearly eliminated and the remaining compromise is
    between stronger pinning and larger photon scattering.
\end{enumerate}

\subsection{Smoothing the cost function}

The raw cost landscape contains narrow oscillatory structures along the power
axis, which are visible in the unsmoothed optimization maps and are naturally
interpreted as features that are sensitive to slow laser-power fluctuations or
power-calibration uncertainty. In the plotting code, this effect is modeled by
applying a one-dimensional Gaussian convolution to the cost function along the
power axis only:
\begin{align*}
    \widetilde{E}_{\rm total}(\lambda,P_i;\sigma_G)
    = \sum_j G_{\sigma_G}(i-j)\,E_{\rm total}(\lambda,P_j),
\end{align*}
where $G_{\sigma_G}$ is a normalized discrete Gaussian with width
$\sigma_G$ measured in grid points. Since both the AC Stark shift and the
scattering rate scale linearly with the local addressing power, this smoothing
can be reinterpreted as an effective quasi-static power noise.

For the present scans, the power grid step is
\begin{align*}
    \Delta P_{\rm grid}
    = \frac{500-140}{40-1}\,\mu\mathrm{W}
    = 9.23~\mu\mathrm{W}.
\end{align*}
Therefore the Gaussian width in the smoothing routine maps to an equivalent
power noise through
\begin{align*}
    \sigma_P = \sigma_G \Delta P_{\rm grid},
    \qquad
    \frac{\sigma_I}{I} \approx \frac{\sigma_P}{P_*},
\end{align*}
where $P_*$ is the operating power around the selected optimum. For example,
$\sigma_G=1$ corresponds to $\sigma_P=9.23~\mu\mathrm{W}$ and
$\sigma_G=3$ corresponds to $\sigma_P=27.69~\mu\mathrm{W}$. At
$P_*=200~\mu\mathrm{W}$ these imply fractional power noise of $4.6\%$ and
$13.8\%$, whereas at $P_*=450~\mu\mathrm{W}$ they imply $2.1\%$ and $6.2\%$,
respectively.

\begin{table}[h]
    \centering
    \caption{Representative smoothed optima for three values of the Gaussian
    width. Each entry lists $(\lambda_{\rm opt},P_{\rm opt})$ in units of
    $(\mathrm{nm},\mu\mathrm{W})$.}
    \label{tab:smoothing_summary}
    \begin{tabular}{lccc}
        \toprule
        Scan & $\sigma_G=0$ & $\sigma_G=1$ & $\sigma_G=3$ \\
        \midrule
        exp1 & $(785.58,\ 204.6)$ & $(784.14,\ 140.0)$ & $(784.14,\ 140.0)$ \\
        exp2 & $(786.00,\ 204.6)$ & $(783.97,\ 140.0)$ & $(784.90,\ 140.0)$ \\
        exp3 & $(785.58,\ 204.6)$ & $(784.14,\ 140.0)$ & $(784.14,\ 140.0)$ \\
        exp4 & $(785.58,\ 490.8)$ & $(785.58,\ 490.8)$ & $(785.32,\ 490.8)$ \\
        exp5 & $(785.24,\ 444.6)$ & $(784.56,\ 306.2)$ & $(783.88,\ 361.5)$ \\
        \bottomrule
    \end{tabular}
\end{table}

Several trends are clear from Table~\ref{tab:smoothing_summary}. First, in
exp1 and exp3 the optimum is quite fragile with respect to power smoothing:
already at $\sigma_G=1$, the optimum moves from the sharp minimum near
$785.6~\mathrm{nm}$ and $205~\mu\mathrm{W}$ to the broader low-power edge near
$784.1~\mathrm{nm}$ and $140~\mu\mathrm{W}$. Second, exp4 remains pinned to the
high-power edge even after smoothing, which confirms that its main limitation
is the too narrow chirp window rather than fine interference structure in the
power direction. Third, exp5 retains the most useful broad minimum: as the
smoothing is increased from $\sigma_G=0$ to $\sigma_G=3$, the optimum moves
continuously from $(785.24~\mathrm{nm},444.6~\mu\mathrm{W})$ to
$(783.88~\mathrm{nm},361.5~\mu\mathrm{W})$, while the minimum cost only rises
from $9.79\times 10^{-4}$ to $1.16\times 10^{-3}$. This means that the
strong-pinning protocol of exp5 is much more robust against power noise than
the shorter and weaker protocols of exp1--exp3.

For exp5, which provides the best strong-pinning setting, the detailed
smoothed optima are particularly informative:
\begin{align*}
    \sigma_G = 0.5 &\Rightarrow (\lambda_{\rm opt},P_{\rm opt})
    = (785.58~\mathrm{nm},500.0~\mu\mathrm{W}),\\
    \sigma_G = 1.0 &\Rightarrow (\lambda_{\rm opt},P_{\rm opt})
    = (784.56~\mathrm{nm},306.2~\mu\mathrm{W}),\\
    \sigma_G = 2.0 &\Rightarrow (\lambda_{\rm opt},P_{\rm opt})
    = (783.88~\mathrm{nm},361.5~\mu\mathrm{W}),\\
    \sigma_G = 3.0 &\Rightarrow (\lambda_{\rm opt},P_{\rm opt})
    = (783.88~\mathrm{nm},361.5~\mu\mathrm{W}).
\end{align*}
Thus the smoothing primarily removes narrow power-dependent fringes and
replaces the raw optimum by a broader valley centered roughly at
$\lambda \approx 784.0$--$784.6~\mathrm{nm}$ and
$P \approx 300$--$360~\mu\mathrm{W}$ once the effective power noise reaches the
few-percent level.

To decide which smoothing level is physically reasonable, we compare these
numbers with the baseline local-laser noise used elsewhere in the simulation,
where a relative intensity noise of $1\%$ is taken as the nominal scale. Near
the strong-pinning exp5 optimum, this corresponds to
$\sigma_P \sim 4$--$5~\mu\mathrm{W}$, or $\sigma_G \sim 0.5$. Therefore,
the default plotting choice $\sigma_G=3$ should be interpreted as a
conservative robustness filter rather than a literal experimental noise model:
it corresponds to several-percent to ten-percent power fluctuations, depending
on the operating power.

Our final smoothing-based decision is therefore twofold. If the local
addressing laser can be stabilized at the $\sim 1\%$ level, the preferred
strong-pinning working point remains close to the raw exp5 optimum, namely
$\lambda \approx 785.5~\mathrm{nm}$ and
$P \approx 450$--$500~\mu\mathrm{W}$. If instead one wants a more conservative
choice that is robust against a few-percent slow power drift, the smoothed
minimum around $\lambda \approx 784.6~\mathrm{nm}$ and
$P \approx 300~\mu\mathrm{W}$ is the better operating point. In this sense,
the smoothing analysis does not invalidate the preference for the
$785~\mathrm{nm}$ region, but it does push the robust power choice downward
once laser noise is taken into account.

\subsection{Recommended wavelength and power}

Based on the five scans above, the wavelength can be recommended more
confidently than the power.

\begin{enumerate}[label=(\arabic*)]
    \item The recommended wavelength window is
    $\lambda \approx 785.3$--$785.8~\mathrm{nm}$, with
    $\lambda \approx 785.5~\mathrm{nm}$ as the central working point. This
    range appears repeatedly in all optimized scans and is therefore more
    robust than the single calibration point at $784~\mathrm{nm}$.

    \item If the goal is to keep the scattering loss low while still obtaining
    a clear local detuning, a practical low-loss operating point is
    $P \approx 180$--$220~\mu\mathrm{W}$. This is the regime selected by exp1,
    exp2, and exp3, and it already yields
    $E_{\rm total}\sim 10^{-4}$--$10^{-3}$.

    \item If the goal is strong pinning with a wider chirp window and a longer
    gate time, exp5 remains the most useful reference scan. Without smoothing,
    its best point is
    $(\lambda,P)\approx(785.24~\mathrm{nm},444.6~\mu\mathrm{W})$. After
    including a moderate smoothing that corresponds to a few-percent power
    fluctuation, the robust optimum moves to about
    $(784.56~\mathrm{nm},306.2~\mu\mathrm{W})$.

    \item Therefore, the most useful practical strategy is to first choose the
    wavelength in the $784.5$--$785.5~\mathrm{nm}$ window and then tune the
    power according to the target protocol and the expected laser stability:
    around $200~\mu\mathrm{W}$ for a low-scattering baseline, around
    $450$--$500~\mu\mathrm{W}$ when stronger pinning is required and the local
    power noise is kept near the $1\%$ level, and around
    $300$--$360~\mu\mathrm{W}$ when a more conservative smoothed optimum is
    preferred.
\end{enumerate}

\subsection{Possible problem}

The present analysis is already sufficient to identify a preferred wavelength
window, but several limitations should be kept in mind.

\begin{enumerate}[label=(\alph*)]
    \item The simulation is based on a two-atom model. It captures the basic
    competition between local pinning, crosstalk, and scattering, but it does
    not yet include full many-body boundary effects, spatially varying local
    weights, or the domain-wall geometry used in the experiment.

    \item The power optimum is protocol dependent. The current data only sample
    a few discrete choices of $(\Delta_{\max},T)$, so they do not yet define a
    unique ``best power'' independent of the sweep protocol. This is why the
    preferred power changes from about $200~\mu\mathrm{W}$ in exp1--exp3 to
    about $450~\mu\mathrm{W}$ in exp5.

    \item In most of the scan, the scattering contribution is added as the
    penalty $1-e^{-\Gamma_{\rm sc}T}$ rather than being included in the full
    dissipative dynamics at every grid point. This approximation is useful for
    fast optimization, but it may underestimate state-dependent decoherence for
    the best candidates.

    \item The nominal adiabatic guideline
    $|\delta_A| \gtrsim 4\Omega_{\rm eff}$ should only be regarded as a rough
    design rule. In the present simulations, the numerical optimum is
    determined by the full balance of leakage, crosstalk, and scattering, and
    it does not always coincide with this simple estimate.

    \item The next step is therefore clear: fix the wavelength near
    $785.5~\mathrm{nm}$, perform a dedicated two-dimensional sweep of
    $(\Delta_{\max},T)$, and then re-evaluate the best few points with
    scattering included directly in the Hamiltonian evolution. After that, the
    same analysis should be extended from the two-atom model to a realistic
    multi-atom geometry with the boundary-weight pattern used in the
    coarsening experiment.
\end{enumerate}

\bibliography{ref}

\appendix


\end{document}
